\scanpage{96}

en faisant attention que pour $i$ et $j$ incomparab\unsure{l} et $X$ metrisable, on \deleted{\unsure{trouve}} aura $V_{ij}=\phi$ \deleted{et} (donc pas la peine d'en parler!), et pour $i=j$, on aura $V_{ij}=V_{ii}=V_i$ (donc pas la peine de l'introduire, $V_i$ suffit!), enfin \deleted{\unsure{et}} si $i$ et $j$ sont distincts et comparables, \deleted{\unsure{d|donc}} comme $V_{ij}=V_{ji}$, il suffit d'introduir $V_{ij}$ pour \deleted{\unsure{les}} $i<j$. Plus généralement, on introduit les
\[(98)\qquad V_{i_0\mathord{-}i_n}=V_{i_0}\cap\mathord{-}\mathord{-}\cap V_{i_n},\]
mais on se \unsure{borne|bornera} en règle générale au cas où $\{i_0,\mathord{-},i_n\}$ est un drapeau (sinon on trouve $\phi$\} et où en fait ils sont rangés en ordre str. croissant
\[i_0<i_1<\mathord{-}\mathord{-}<i_n\]
\unsure{(}pour éviter des redondances). Les protopos $^*$ sont les ``pierres de construction'' élémentaires, pour en construire des tas d'autres, par $\varinjlim$ (finies si $I$ fini), compte tenu des morphismes semi-simpliciaux évidents entre les $V_d$ $(d=(i_0,\mathord{-},i_n)\in\mathrm{Drap}(I))$.
\margin{$^*$(et des protopos homotopes $V_{i_0,\mathord{-}i_n}^{Z}$ plus bas)}
\margin{95}
